% !TeX root = ../../main.tex

Der Codeausschnitt~\ref{lst:params} zeigt die Wertebereiche der für die Experimente verwendeten Parameter.
\texttt{WORKLOAD} und \texttt{ITERATION} bleiben über alle Experimente hinweg konstant und umfassen jeweils
500 auszuführende Transaktionen sowie fünf Wiederholungen.
Dadurch wird verhindert, dass Unterschiede in der Anzahl der ausgeführten Transaktionen die Messergebnisse
beeinflussen.
Unterschiede zwischen den Ergebnissen können somit auf die jeweils untersuchten Parameter zurückgeführt werden.

Die Parameter \texttt{CONCURRENCY}, \texttt{RETRY\_COUNT} und \texttt{RETRY\_DELAY} werden jeweils über mehrere
festgelegte Werte variiert.
Dadurch kann untersucht werden, wie sich beispielsweise eine steigende Anzahl paralleler Clients, eine höhere
Anzahl zulässiger Retries oder ein zunehmender Delay auf die jeweiligen Kenngrößen auswirkt.
Die Experimente werden dabei so konfiguriert, dass der jeweils untersuchte Parameter variiert wird, während die
übrigen Parameter konstant gehalten werden.
Dadurch lässt sich der Einfluss des einzelnen Parameters gezielter bestimmen.

Der Parameter \texttt{RETRY\_STRATEGY} ist aus Gründen der Optimierung ebenfalls in der \texttt{config.py} definiert.
Die Auswahl der tatsächlich verwendeten Strategie erfolgt jedoch innerhalb der Implementierung über eine
\texttt{match/case}-Verzweigung.
Dadurch können die verschiedenen Retry-Strategien mit derselben grundlegenden Logik ausgeführt und miteinander
verglichen werden.

Die gewählten Parameterbereiche stellen dabei einen Kompromiss zwischen dem Umfang der Experimente und der
Aussagekraft der Ergebnisse dar.
Durch die Variation mehrerer Werte können Trends innerhalb der untersuchten Kenngrößen erkannt werden, ohne dass
für jede mögliche Kombination aller Parameter ein eigenes Experiment durchgeführt werden muss.
Die fünf Wiederholungen pro Konfiguration dienen zusätzlich dazu, Schwankungen einzelner Messungen zu reduzieren
und eine stabilere Grundlage für die spätere Auswertung zu schaffen.

\begin{lstlisting}[style=pythoncode, caption={Übersicht der Parameter}, label={lst:params}]
RETRY_COUNT = [0, 1, 3, 5]
RETRY_DELAY = [0, 0.005, 0.01, 0.05, 0.1]
RETRY_STRATEGY = [0, 1, 2, 3]

CONCURRENCY = [2, 5, 7, 10, 20]
WORKLOAD = [500]
ITERATIONS = 5
\end{lstlisting}